% wave15_n5_HH_E3_compact_vs_open.tex
% Wave-15 N5/20: $\HH^*_{E_3}(\Obs_{\hCS})$ on non-compact $\C^3$ vs compact CY$_3$.
% Channels: Lurie + Gwilliam. Five attack-heal cycles.
% Primary: Lurie 2009; Gwilliam--Williams 2021; Ayala--Francis 2015; Costello 2013.

\providecommand{\Obs}{\mathrm{Obs}}
\providecommand{\hCS}{\mathrm{hCS}}
\providecommand{\CE}{\mathrm{CE}}
\providecommand{\Ch}{\mathrm{Ch}}
\providecommand{\Dolb}{\mathrm{Dolb}}
\providecommand{\chir}{\mathrm{chir}}
\providecommand{\dbar}{\bar\partial}
\providecommand{\kch}{\kappa_{\mathrm{ch}}}
\providecommand{\kcat}{\kappa_{\mathrm{cat}}}
\providecommand{\HH}{\mathrm{HH}}
\providecommand{\Sym}{\mathrm{Sym}}
\providecommand{\Tr}{\mathrm{Tr}}
\providecommand{\C}{\mathbb{C}}
\providecommand{\cO}{\mathcal{O}}
\providecommand{\cE}{\mathcal{E}}
\providecommand{\g}{\mathfrak{g}}
\providecommand{\fgl}{\mathfrak{gl}}
\providecommand{\Perf}{\mathrm{Perf}}
\providecommand{\Mod}{\mathrm{Mod}}
\providecommand{\Alg}{\mathrm{Alg}}
\providecommand{\Bord}{\mathrm{Bord}}
\providecommand{\Prim}{\mathrm{Prim}}

\section*{$\HH^*_{E_3}(\Obs_{\hCS})$: non-compact $\C^3$ vs compact CY$_3$ (Wave-15 N5/20)}

\paragraph{Theorem (scope comparison).}
The $E_3$-Hochschild cohomology of $\Obs_{\hCS}$ separates sharply by
ambient geometry. On $\C^3$, $\HH^0_{E_3}(\Obs_{\hCS}(\C^3, \g))$ is
an infinite-rank formal power series algebra
$\C[\![\tau_1, \tau_2, \tau_3]\!]$ (Gwilliam--Williams 2021 Prop.~5.3.2);
on compact CY$_3$ it is finite-rank with rank computed by Hodge data
and the Chevalley--Eilenberg cohomology of $\g$. The failure of
non-abelian 3-dualizability on $\C^3$ (Wave-14 J2) therefore recovers
on compact CY$_3$; the cobordism hypothesis (Lurie 2009 Thm~2.4.6)
upgrades 6D hCS to a fully extended framed 3-TFT only after
compactification. At $d=3$ the CY-to-chiral functor $\Phi$ is
$E_1$-chiral (the $E_2$-braiding lives on $Z(\mathrm{Rep}(A))$,
not on $A$); compactness supplies precisely the perfectness
required for extended functoriality.

\paragraph{Cycle 1: $\HH^0_{E_3}$ on $\C^3$ is $\C[\![\tau_1, \tau_2, \tau_3]\!]$.}
\emph{Attack.} Francis 2013 Thm~1.1 identifies $\HH^*_{E_n}(A)$ with
the cotangent complex of $A$ as an $E_n$-algebra: morally
$\HH^*_{E_n}(A) \simeq \mathrm{Map}_{\Alg_{E_n}}(A, A)$, concentrated
by the Poincar\'e--Birkhoff--Witt filtration into symmetric-algebra
corrections of the primitives $\Prim(A)$.
\emph{Heal.} For $A = \Obs_{\hCS}(\C^3, \fgl_n)$, the PBW presentation
$A \simeq U_{E_3}(\g_{\hCS})$ (Wave-13 H4) with
$\g_{\hCS} = \Omega^{0,\bullet}(\C^3, \fgl_n)[1]$ gives
\[
  \HH^0_{E_3}(A) \;\simeq\; \HH^*_{\mathrm{Lie}}\bigl(\g_{\hCS}, \cO(\C^3)\bigr)^{E_3}
  \;=\; \bigl(\Sym^\bullet \Omega^{0,0}(\C^3)\bigr)^{\fgl_n}
  \;\simeq\; \C[\![\tau_1, \tau_2, \tau_3]\!],
\]
the three Casimir generators $\tau_i = \mathrm{Tr}(z_i^{\partial}) \in
\Omega^{0,0}(\C^3)^{\fgl_n}$ parametrising the three translation
directions on $\C^3$ (Gwilliam--Williams 2021 Prop.~5.3.2, stated
there for the higher Kac--Moody algebra $\widehat{\fgl}_n^{(3)}$).
Infinite-dimensionality is forced by the fact that $\cO(\C^3)$ is a
polynomial ring in three variables and the Casimir trace descends
to each polynomial monomial.

\paragraph{Cycle 2: finiteness on compact $X$ from Hodge truncation.}
\emph{Attack.} Replace $\C^3$ by a compact CY$_3$ $X$. The chiral
CE complex becomes $\CE^*_{\dbar,\chir}(\Omega^{0,\bullet}(X, \g), \cO_X)$;
pushforward to a point is given by $\dbar$-cohomology of $X$,
i.e. the Dolbeault complex $\bigoplus_q H^{0,q}(X)$.
\emph{Heal.} The finite-rank statement is: $\HH^0_{E_3}(\Obs_{\hCS}(X, \g))
\simeq H^0(X, \Sym^\bullet \g^\vee)^{\g}$ whose bigraded total rank is
\[
  \dim \HH^0_{E_3}(\Obs_{\hCS}(X, \g)) \;=\;
  \sum_{q \geq 0} h^{0,q}(X) \cdot \dim \HH^*_{\mathrm{Lie}}(\g, \C)^{[q]} \;<\; \infty,
\]
because $h^{0,q}(X)$ vanishes outside $q \in \{0, 1, 2, 3\}$ on a
compact threefold and $\HH^*_{\mathrm{Lie}}(\g, \C)$ is finite for
reductive $\g$ (Chevalley--Eilenberg). Coherence of $\Obs_{\hCS}$ as
a sheaf of $E_3$-algebras (Costello 2013 \S8 / Costello--Gwilliam
Vol~II Thm~9.3.1) plus compactness of $X$ gives finiteness of
pushforward to a point (Grothendieck finiteness).

\paragraph{Cycle 3: $K3\times E$ via K\"unneth and classical $\HH^*$.}
\emph{Attack.} On $X = K3 \times E$, the $E_3$-Hochschild complex
factorises by the holomorphic K\"unneth formula
$\CE^*_{\dbar,\chir}(\Omega^{0,\bullet}(X, \g)) \simeq
\CE^*_{\dbar,\chir}(\Omega^{0,\bullet}(K3, \g)) \otimes
\CE^*_{\dbar,\chir}(\Omega^{0,\bullet}(E, \g))$, mirroring the
multiplicativity of $\kcat$ under products.
\emph{Heal.} The Hodge data $h^{0,0}(K3) = h^{0,2}(K3) = 1$,
$h^{0,1}(K3) = 0$, and $h^{0,0}(E) = h^{0,1}(E) = 1$ give
\[
  \HH^0_{E_3}(\Obs_{\hCS}(K3 \times E, \g)) \;\simeq\;
  \bigl(\C \oplus \C[-2]\bigr) \otimes \bigl(\C \oplus \C[-1]\bigr) \otimes
  \HH^*_{\mathrm{Lie}}(\g, \C)^{E_3},
\]
of finite total rank $4 \cdot \dim \HH^*_{\mathrm{Lie}}(\g, \C)^{E_3}$.
This is consistent with
$\kcat(K3 \times E) = \chi(\cO_{K3 \times E}) = \chi(\cO_{K3}) \cdot
\chi(\cO_E) = 2 \cdot 0 = 0$ (K\"unneth-multiplicative on the total
space; the value $2$ is the K3 fibre alone, not the product), and
with $\kch(\cA_{K3 \times E}) = 0$ via the Hodge-supertrace
identification (cy\_d\_kappa\_stratification.tex, line 755). The
finite rank of $\HH^0_{E_3}$ matches the Mukai-enhanced Heisenberg
ceiling $K^{\kappa_{\mathrm{ch}}} = 8$ on the $\mathsf B$-row of Theorem C through
the $K3\times E$ specialisation.

\paragraph{Cycle 4: dualizability consequence.}
\emph{Attack.} Lurie HA \S5.2 / Francis 2013 Thm~3.4: an $E_n$-algebra
$A$ is $n$-dualizable iff $A$ is perfect over itself, equivalently iff
$\HH^*_{E_n}(A)$ is dualizable in $\Ch$.
\emph{Heal.} On $\C^3$: $\C[\![\tau_1, \tau_2, \tau_3]\!]$ is not
dualizable in $\Ch$ (an infinite-rank formal power series has no
finite dual), so $\Obs_{\hCS}(\C^3, \g)$ fails 3-dualizability
non-abelianly (Wave-14 J2 Cycle 4); 6D hCS on $\C^3$ extends only
to a $(0,1,2)$-truncated framed $E_3$-TFT. On compact CY$_3$:
the finite rank of Cycle 2 makes $\HH^*_{E_3}$ dualizable, so
$\Obs_{\hCS}(X, \g)$ is 3-dualizable, and the Calaque--Pantev--
Toën--Vaquié--Vezzosi 2017 Prop.~2.6 PTVV $(-3)$-shifted symplectic
form supplies the evaluation/coevaluation data. Lurie 2009 Thm~2.4.6
then upgrades 6D hCS on $X$ to a fully extended framed 3-TFT.

\paragraph{Cycle 5: link to CY-D $\Phi$-output scope.}
\emph{Attack.} CY-D stratifies $\Phi: \mathrm{CY}\text{-cat}_d \to
\mathrm{ChirAlg}$ by output type: $E_2$-chiral at $d \leq 2$,
$E_1$-chiral at $d \geq 3$. At $d = 3$ the output is $E_1$;
$E_2$-braiding lives on $Z(\mathrm{Rep}(A))$, not on $A$.
\emph{Heal.} The Cycle 4 dualizability gap tracks this stratification
precisely. On $\C^3$, non-compactness obstructs both $E_3$-perfectness
(via the infinite $\HH^0_{E_3}$) and the extended functoriality of
$\Phi$ as an $(\infty, 3)$-categorical morphism; the object-level
$\Phi$ still produces an $E_1$-chiral algebra on each transverse
curve $C \subset \C^3$, but the 3-morphism duals that would upgrade
$\Phi$ to a fully extended functor require compactness. On compact
$X$, the finite $\HH^*_{E_3}$ supplies the 3-morphism duals;
$\Phi$ extends to an $(\infty,3)$-functor on the subcategory of
compact CY$_3$ inputs. The ``$\Phi$ is $E_1$-chiral at $d\geq 3$''
statement is the chain-level shadow of this: once you drop from
$E_3$ down to $E_1$ at the output, the functoriality loss on
non-compact inputs is masked; the extended-functoriality gap is
genuinely an $E_3$-level phenomenon, visible only when one asks
about 2-morphism and 3-morphism duals.

\paragraph{Relation to $\kch$, $\kcat$ and BKM.}
The trace $\tau$ of Wave-14 J2 Cycle 1 realises the $\Phi$-side
chiral Euler character
$\kch(\cA_{\hCS, X}) = \sum_q (-1)^q h^{0,q}(X)$ on compact CY$_3$.
Finite $\HH^*_{E_3}$ is the chain-level witness that this supertrace
converges. On $\C^3$ the supertrace is ill-defined (infinite
Fock-space trace); on compact $X$ the Hodge truncation makes it a
well-defined integer, matching $\kcat = \chi(\cO_X)$ under the
CY-D stratification. The Borcherds weight identity
$\kappa_{\mathrm{BKM}}(\Phi_N) = c_N(0)/2$ is the
automorphic-correction side and is independent of this
$\HH^*_{E_3}$ finiteness question.

\paragraph{Scope and lanes.}
Chain-level: explicit $\C[\![\tau_1,\tau_2,\tau_3]\!]$ presentation
on $\C^3$; explicit K\"unneth decomposition on $K3 \times E$;
explicit Hodge truncation on compact $X$. $(\infty,1)$-categorical:
Lurie HA \S5.2 dualizability; Francis 2013 Thm~3.4
perfectness criterion; Lurie 2009 cobordism hypothesis;
Calaque--Pantev--Toën--Vaquié--Vezzosi 2017 Prop.~2.6
$(-n)$-shifted symplectic $\Rightarrow$ $n$-dualizable. Both lanes
load-bearing (CLAUDE.md equal-status rule).

\paragraph{References.}
\begin{itemize}
  \item J.\ Lurie, ``On the classification of topological field
    theories,'' 2009, Thm~2.4.6; \emph{Higher Algebra} \S5.2--5.3.
  \item O.\ Gwilliam and B.\ R.\ Williams, ``Higher Kac--Moody
    algebras and symmetries of holomorphic field theories,''
    arXiv:2009.05037 (2021), Prop.~5.3.2.
  \item D.\ Ayala and J.\ Francis, ``Factorization homology of
    topological manifolds,'' \emph{J.\ Topology} 8 (2015), Thm~1.1.
  \item J.\ Francis, ``The tangent complex and Hochschild cohomology
    of $E_n$-rings,'' \emph{Compos.\ Math.} 149 (2013), Thm~1.1, 3.4.
  \item K.\ Costello, ``Notes on supersymmetric and holomorphic field
    theories,'' arXiv:1111.4234 (2013), \S8.
  \item K.\ Costello and O.\ Gwilliam, \emph{Factorization Algebras
    in QFT} Vol~II, Thm~9.3.1.
  \item D.\ Calaque, T.\ Pantev, B.\ Toën, M.\ Vaqui\'e, G.\ Vezzosi,
    arXiv:1506.03699, Prop.~2.6.
\end{itemize}
